{\scriptsize
    \begin{minipage}[t]{0.5\textwidth}
\begin{tabularx}{\linewidth}{|p{0.6\textwidth}|X|}
\hline
Notions et contenus & Capacités exigibles \\
\hline
1.3. Exemple de dispositif interférentiel par division du front d'onde : trous d'Young & \\
\hline
Dispositif-modèle des trous d'Young ponctuels dans un milieu non dispersif (source ponctuelle à grande distance finie ; observation à grande distance finie).
Champ d'interférences. Ordre d'interférences. & Définir, déterminer et utiliser l'ordre d'interférences. \\
\hline
Franges d'interférences. & Justifier la forme des franges observées sur un écran éloigné parallèle au plan contenant les trous d'Young. \\
\hline
Du dispositif-modèle au dispositif réel. & \\
\hline
Fentes d'Young.
Expliquer l'intérêt pratique du dispositif des fentes d'Young comparativement aux trous d'Young.  &  Identifier l'effet de la diffraction sur la figure observée. Montage de Fraunhofer.  \\
\hline

\end{tabularx}
\end{minipage}%
\begin{minipage}[t]{0.5\textwidth}
    \begin{tabularx}{\linewidth}{|X|X|}
\hline 
Exprimer l'ordre d'interférences sur l'écran dans le cas d'un dispositif des fentes d'Young utilisé en configuration de Fraunhofer. & \\
\hline
Perte de contraste par élargissement spatial de la source. & Utiliser un critère semi-quantitatif de brouillage des franges portant sur l'ordre d'interférences pour interpréter des observations expérimentales. \\
\hline
Perte de contraste par élargissement spectral de la source. & Utiliser un critère semi-quantitatif de brouillage des franges portant sur l'ordre spectral de la source. Relier la longueur de cohérence temporelle, la largeur spectrale et la longueur d'onde en ordres de grandeur. \\
\hline
Observations en lumière blanche (blanc d'ordre supérieur, spectre cannelé). & Déterminer les longueurs d'ondes des cannelures. \\
\hline
\end{tabularx}
\end{minipage}
}
